\subsection{Effective separably closed fields of higher imperfection degree}
\label{C:separably-closed-subsection}

\begin{theorem}[Proposed: exact separable closures and finite-system decisions]
\label{C:separably-closed-effective-arithmetic}
For every odd prime power $q$ of characteristic $p$
and every finite $r\ge1$, there is an explicitly
constructed countable presentation
\[
       S_r=\mathbf F_q(t_1,\ldots,t_r)^{\rm sep}        \tag{C.3.SC1}
\]
with decidable exact finite names, equality and
arithmetic.  The elements $t_1,\ldots,t_r$
form a $p$-basis, and their $p$-component
functions are computable.
Every finite system of polynomial equations
and inequations over $S_r$ has an ordinary
terminating solution decision; a positive
answer constructs an exact solution.

Uniformly, the same solution decisions hold
for finite-dimensional commutative $S_r$-algebras
given by tables, by writing variables in their
supplied $S_r$-bases.
All finite field towers over $S_r$ have
decidable validity, exact arithmetic and
ordinary univariate factorization, including
inseparable towers and polynomials.
Every finite field extension of $S_r$ is
purely inseparable and separably closed.
No unspecified equation-solving oracle is
assumed in these assertions.
\end{theorem}
\begin{proof}
Put $k_0=\mathbf F_q(t_1,\ldots,t_r)$.
First there is ordinary factorization over
$k_0$.  Clear denominators of a univariate
polynomial and remove its content to get
a primitive polynomial
$F\in\mathbf F_q[t_1,\ldots,t_r,X]$.
If its total degree is $D$, every polynomial
factor has total degree at most $D$.
There are only finitely many polynomials
over $\mathbf F_q$ of that degree bound.
Gauss' lemma says that every factorization
over $k_0[X]$ is detected by exact
polynomial division among these candidates.
Normalize the factors and recover their
multiplicities by division.
Zero polynomials, constants, and leading
coefficients are handled separately.
This is finite enumeration, with no
complexity claim.

Separable polynomials can then be factored
over any explicit finite separable tower
over $k_0$.  The separable-extension part
of Theorem~\ref{C:function-field-factorization}
applies to any infinite effective ground
field with factorization for separable
polynomials.  Specifically, for
$L=E(\alpha)$ finite separable and monic
separable $f\in L[X]$, enumerate $c\in E$
until
\[
            N_{L/E}(f(X-c\alpha))                     \tag{C.3.SC2}
\]
is separable.  The norm is a determinant.
Distinct conjugate roots collide for only
finitely many $c$, so the search terminates.
Factor (C.3.SC2) over $E$ and take its
gcds with $f(X-c\alpha)$ over $L$.
Separability of the norm makes these
exactly the irreducible factors, and
the shift can be undone.
Only separable input is required in
the tower construction below.

Build an increasing sequence of finite
separable fields $k_n/k_0$ with explicit
embeddings.  At step $n$, consider all
nonconstant polynomial codes at most $n$ over all
already constructed stages $k_i$, $i\le n$.
Discard invalid or nonseparable polynomials.
Using the preceding factorization, adjoin
enough roots to split the remaining finite
list over the current field.
Every adjoined root is separable and
the list requires finitely many adjunctions.
This defines the next explicit finite
separable stage $k_{n+1}$.
Every separable polynomial over the union
has its coefficients at a finite stage
and occurs in this fair schedule.
Hence the union is separably closed and
separable algebraic over $k_0$; it is a
presentation of (C.3.SC1).
A name consists of a stage number and
a finite field-coordinate expression.
Validity is decidable by constructing
that finite stage; arithmetic and equality
are computed after transport to a common
later stage.

We verify and compute the $p$-basis.
For any finite separable $L/k_0$ of
degree $d$, the extension
$L/k_0L^p$ is both purely inseparable
and separable, hence trivial.
Frobenius identifies $L/k_0$ with
$L^p/k_0^p$, so
\[
 [L:L^p]=p^r,
 \qquad
 \{t_1^{a_1}\cdots t_r^{a_r}:0\le a_i<p\}
        \text{ is an }L^p\text{-basis of }L.           \tag{C.3.SC3}
\]
The monomials span because $L=k_0L^p$,
and the displayed degree makes them
independent.
Passing to the union preserves spanning;
any putative dependence occurs already
at one finite stage and is therefore
zero.  This proves that the same
monomials are an $S_r^p$-basis.

For an explicit $x\in L=k_n$, choose
the known $k_0$-basis $e_1,\ldots,e_d$.
The elements $t^a e_i$ form a
$k_0^p$-basis of $L$.
The elements $t^a e_i^p$ form another
such basis, by (C.3.SC3) and the
Frobenius image of the $k_0$-basis.
Their change-of-basis matrix is
computable over
$k_0^p=\mathbf F_q(t_1^p,\ldots,t_r^p)$.
For completeness, a rational function
$A/B\in k_0$ is expressed in its
$p$-monomial basis by writing it as
$AB^{p-1}/B^p$ and grouping numerator
exponents modulo $p$.
Solving this finite linear system
expresses $x$ uniquely as
\[
              x=\sum_a t^a c_a^p,
                     \qquad c_a\in L.                 \tag{C.3.SC4}
\]
The required $p$th roots of coefficients
in $k_0^p$ are computed by dividing
exponents and applying inverse Frobenius
in the finite field.
Thus each component $c_a$ in (C.3.SC4)
has a computed exact finite name.

Haran's explicit quantifier-elimination
theorem applies in precisely this
language: field operations, a named
finite $p$-basis, and the finitely
many component functions in (C.3.SC4)
\cite[Theorem 1.2 and Sections 1--3,
pp. 464--469]{C-HaranSeparablyClosed1988}.
Given a finite polynomial system,
regard its named coefficients as
free parameters and eliminate its
quantified unknowns by that algorithm.
The resulting quantifier-free formula
is evaluated using the arithmetic,
equality and component algorithms
just proved.  Inequations can be
retained in the formula, or expressed
by adjoining inverses.
If a solution exists, enumeration
of exact tuples and the decidable
original equations finds one.
This proves the asserted terminating
decision and witness construction.
The cited theorem supplies an actual
effective syntactic procedure, not
a hypothesis about an external
existential oracle.

In a finite-dimensional commutative
$S_r$-algebra $T$, replace each
unknown by its coordinates in the
given basis.  Its multiplication
table turns a polynomial equation
into finitely many polynomial
equations over $S_r$.
An inequation is the finite
disjunction that one coordinate
is nonzero.  Haran's decision
therefore also handles every
finite system over $T$.

For a monic degree-$d$ polynomial
over a supplied finite field $L/S_r$,
proper monic factors of degrees
$j,d-j$, $1\le j<d$, are unknown
finite coordinate vectors over
$S_r$.  Their product equals the
input exactly when an explicit
finite polynomial system holds.
The preceding decision tests all
these finitely many possibilities.
If none exists the input is
irreducible; otherwise positive
enumeration finds a factor pair.
Recurse by degree and restore
the leading scalar.
This factors all polynomials,
with zero and constant cases
treated directly.
Inductively the same procedure
validates each defining polynomial
before admitting the next tower
field.  There is no circular use
of field validity.

Finally every finite extension
of a separably closed field is
purely inseparable: the maximal
separable subextension is trivial
\cite[Lemma 9.14.6, Tag 09HD]
{C-StacksSeparablyClosed}.
An algebraic closure of $S_r$
is purely inseparable over it,
and therefore purely inseparable
over any intermediate finite
extension.  That extension is
again separably closed.
\end{proof}

\begin{theorem}[Proposed: canonical free-generator split groups over separably closed coefficient bases]
\label{C:separably-closed-coefficient-WG}
Let $k$ be a supplied finite extension of
$S_r$ from
Theorem~\ref{C:separably-closed-effective-arithmetic}.
For each $m\ge0$, complete ordinary
canonical classification, equality,
product and exterior-operation algorithms
for $WG_{\rm sp}(BO_{m,R})$ hold for:
\begin{enumerate}
\item every finite-dimensional commutative
$k$-algebra $R$ given by its table;
\item every $R=B^h$ with the effective
coherent local input of
Theorem~\ref{C:effective-coherent-henselization-WG},
except that its residue field is now
a supplied finite extension of $k$.
\end{enumerate}
All indecomposable finite projective
rational representations have a literal
decidable canonical index set.
Every self-dual indecomposable has
exactly one nonsingular invariant
form sign.
Let $\mathscr I_+$ and $\mathscr I_-$
be respectively its positive and
negative sign indices, and let
$\mathscr P$ be the unordered
nonself-dual pairs.
There is a canonical free abelian
description
\[
 WG_{\rm sp}(BO_{m,R})\simeq
 \bigoplus_{i\in\mathscr I_+}\mathbf Z[q_i]
 \ \oplus\
 \bigoplus_{i\in\mathscr I_-}\mathbf Z[H(M_i)]
 \ \oplus\
 \bigoplus_{\{i,i^*\}\in\mathscr P}\mathbf Z[H(M_i)] . \tag{C.3.SC5}
\]
Here $q_i$ is a computed reference
symmetric form and $H(M_i)$ is the
actual hyperbolic form.
Every finite coordinate vector has
computed positive and negative
representatives.  All actual isometry
decisions and all finite virtual
operations terminate.
The residue division fields of the
indecomposables are computed finite
purely inseparable extensions of
the ground or residue field; they
are not silently identified with
that field.
\end{theorem}
\begin{proof}
The arithmetic theorem gives every
finite-tower factorization and
finite linear-algebra operation
used in the earlier finite-algebra
and coherent-local pipelines.
We identify why division factors
and multiplicity forms are fully
decided here without assuming $C_1$.

If $L$ is any finite field
extension of $k$, or of the
specified local residue field,
it is finite over $S_r$ and
separably closed.
Every central simple algebra
over $L$ has a finite separable
splitting field
\cite[Proposition 11.8.5, Tag 0752]
{C-StacksGlobalFunctionArithmetic}.
There is no nontrivial such
extension of $L$, so the algebra
is already a matrix algebra.

For an arbitrary finite-dimensional
algebra $E$ over $k$, enumerate
finite field towers $L_j/k$,
matrix sizes, and unital
multiplicative surjections
\[
                 E\longrightarrow\prod_j M_{s_j}(L_j)
\]
whose computed kernel has
$(\dim_kE+1)$st power zero.
Every condition is a finite
decidable test by the arithmetic
theorem.  Such a kernel is
exactly the radical, and
Artinian structure together
with the preceding splitting
argument guarantees a certificate.
Nilpotent lifting gives actual
orthogonal primitive idempotents.
Thus the finite-algebra algorithm
computes all division quotients,
which are the actual fields
$L_j$, including their possibly
nontrivial purely inseparable
extensions of $k$.
The same algorithm applies
over the local residue field.
The zero algebra has empty quotient
product and is handled immediately.

The zero coefficient ring has the
zero representation category and
zero split group.  For the other
finite coefficient algebras,
projectivity is the finite affine
linear section test for the
canonical free-module surjection;
Hom spaces, $R$-duals and invariant
sign forms are actual finite
$k$-linear systems.
Determinant grids over the
infinite field $k$ decide their
invertible members.
The computed endomorphism radical
and primitive lifts give a
terminating actual decomposition.
For $R=B^h$, the exact standard
\'{e}tale-neighborhood kernel proof
of Theorem~\ref{C:effective-coherent-henselization-WG}
does not use the $C_1$ property.
It supplies full actual Hom
kernels and their actual
residue images.  Determinant
grids are applied to these
images, not to all residue
intertwiners.
The kernel of the actual
endomorphism-image map is
radical by the unit determinant
criterion, and
Lemma~\ref{C:nonnoetherian-henselian-lifting}
provides actual idempotent lifts
and form-isometry lifts.
Ranks decrease through every
nontrivial splitting.
Zero modules are handled separately.
Finite predecessor tests on
these decidable data give
the claimed canonical indices
and computed duality, exactly
as in the full proof of
Theorem~\ref{C:henselian-global-coefficient-WG}.

For a self-dual indecomposable
the $g+g^{\mathsf t}$ and
$g-g^{\mathsf t}$ argument
in its local endomorphism
ring produces at least one
nonsingular sign form.
Let $q$ be one, of sign
$\varepsilon$, and let $D=L$
be the actual division quotient.
Its adjoint fixes the ground
field, so it is the identity
on $L$: every element of
this purely inseparable field
extension has a $p$th-power
iterate in the fixed field,
and such a root is unique.
A form of the opposite sign
would give a unit whose residue
satisfies $x=-x$ in $L$.
Since $p>2$, this is impossible.
Thus exactly one sign occurs,
and the determinant tests
compute it and a reference form.

Every nonzero element of a
separably closed field of odd
characteristic is a square,
since $X^2-a$ is separable.
Symmetric field forms therefore
diagonalize to the standard
form and are classified by
their rank.  Alternating field
forms have even rank and
their usual symplectic basis.
All required roots and basis
changes are computable: root
existence is proved and exact
enumeration finds a root.

The actual radical transfer
now gives (C.3.SC5).
On a positive-sign indecomposable,
symmetric multiplicity forms
are unique in each size, with
generator $q_i$.
On a negative-sign indecomposable,
the multiplicity form must be
alternating, so its least
positive size is two.
The resulting symmetric form
is $H(M_i)$: identifying the
dual using the negative reference
form gives the block matrix
$\left(\begin{smallmatrix}0&q_i\\-q_i&0\end{smallmatrix}\right)$.
A nonself-dual pair has the
usual unique hyperbolic block.
The semisimple block multiplicities
are independent, so the monoid
of symmetric forms is free
commutative on these generators;
its group completion is the
displayed free abelian group.

All residual standard matrices
lift through computed actual
endomorphism generators and
are symmetrized by $1/2$.
They are nonsingular modulo
the radical and hence actually
nonsingular.  Radical isometry
lifting, or exact enumeration
with that existence proof,
constructs every desired actual
isometry.
Taking direct sums and formal
differences realizes every
finite integer vector.

Finally tensor products and
exterior powers use the finite
actual projective modules,
balanced tensor quotients and
determinant pairings already
proved in
Theorem~\ref{C:henselian-global-coefficient-WG}.
Reclassify these outputs by
the present terminating block
algorithm.  Bilinearity and
the finite recurrence for
$\lambda_t(x-y)=\lambda_t(x)\lambda_t(y)^{-1}$
compute every requested virtual
coefficient, without dividing
by a factorial.
\end{proof}

\begin{corollary}[Proposed: concrete non-$C_1$ coefficient families]
\label{C:separably-closed-nonC1-family}
For every finite $r\ge2$, the fields
$S_r$ in (C.3.SC1) are separably closed
but are not $C_1$.
Consequently
Theorem~\ref{C:separably-closed-coefficient-WG}
gives complete ordinary split-$WG$
answers for a whole new family of
fields and nonfields, including
$S_r[\epsilon]/(\epsilon^s)$,
$s\ge2$, and $S_r[z]_{(z)}^h$.
More generally it includes all
finitely presented local henselizations
at supplied points with finite
residue extensions over $S_r$,
and all effective coherent local
henselizations with those residue
fields as specified in the theorem.
\end{corollary}
\begin{proof}
The $p$-monomial basis in
(C.3.SC3) gives the homogeneous
polynomial
\[
        \sum_{0\le a_1,\ldots,a_r<p}
                    t^a X_a^p .                       \tag{C.3.SC6}
\]
It has degree $p$ and $p^r>p$
variables.  Its only zero over
$S_r$ is the zero vector,
because the coefficients $t^a$
are linearly independent over
$S_r^p$.
This is an explicit violation
of the $C_1$ condition.
These fields are consequently
not algebraically closed, and
their finite purely inseparable
division-field extensions need
not be trivial.

The field $S_r$ has transcendence
degree $r\ge2$ over the finite
prime field.  Hence it is
neither a global function
field nor a finite extension
of a henselized global
function field, all of which
have transcendence degree one.
It has positive characteristic,
so the $\mathbf Q$-algebra
answer does not apply.

Finite polynomial syzygies over
the effective field $S_r$
and the exact pointed-localization
proof of
Lemma~\ref{C:C1-henselization-access}
verify the kernel hypotheses
for every stated affine local
base; that access proof uses
no $C_1$ argument.
The displayed nonfield examples
are therefore direct instances
of the preceding theorem,
and may be nonreduced or
positive dimensional.
\end{proof}
